% ============================================================
% §2.4 · Diseño Algorítmico — Manual Técnico K-QGMIP
% ============================================================

\section{Diseño Algorítmico}

\subsection{Visión general}

Ambas estrategias comparten la misma filosofía de diseño: \textbf{anclar la solución exacta
de $k=2$ y refinarla iterativamente} hasta obtener $k$ partes, evitando la explosión
combinatoria de la búsqueda exhaustiva.

\begin{itemize}
  \item \textbf{KGeoMIP (E4)}: ancla en GeoMIP (bipartición geométrica exacta vía tabla BFS
        sobre el hipercubo de Hamming), construye una matriz de similitud causal $S$ a partir
        de esa tabla, y en cada paso refina la parte de menor $\Delta\Phi$ usando la heurística
        guiada por $S$.

  \item \textbf{KQNodes (C4)}: parte del conjunto completo $V$ de dimensiones como una sola
        parte, la bipartite con Queyranne (Ordenamiento de Máxima Adyacencia, MAO), e itera
        $k-1$ veces eligiendo siempre la parte de menor $\varphi_{\text{local}}$ para
        bipartirla a continuación.
\end{itemize}

Ambas estrategias calculan la \textbf{EMD-Effect real una sola vez al final}, sobre la
k-partición ganadora. Durante el refinamiento usan funciones proxy más baratas
($\Delta\Phi$ incremental en KGeoMIP, función $f_{\text{local}}$ en KQNodes) que evitan
evaluar la EMD completa en cada paso.

% ─────────────────────────────────────────────────────────────────────────────
\subsection{Estructuras de datos}
% ─────────────────────────────────────────────────────────────────────────────

\subsubsection*{NCube}

Un \textbf{NCube} es el hipercubo de probabilidades de transición de un nodo futuro $i$.
Su atributo central es \texttt{data}: un \texttt{ndarray} de forma $(2, 2, \ldots, 2)$
con $D$ ejes binarios (uno por dimensión presente del subsistema), donde
\texttt{data[$b_{D-1}$, \ldots, $b_0$]} = $P(\text{nodo}_i = \text{ON} \mid s_t = b)$.
La operación fundamental sobre un NCube es \texttt{marginalizar(dims)}: promedia
sobre los ejes de las dimensiones en \texttt{dims}, reduciendo el tensor a los ejes
restantes.

\subsubsection*{Tabla de costos $T$ (GeoMIP)}

La tabla $T \in \mathbb{R}^{2^D \times D_{\text{nc}}}$ es construida por GeoMIP
mediante un recorrido BFS sobre el hipercubo de Hamming:
\begin{itemize}
  \item Fila $s$: estado binario del presente ($s \in \{0,1\}^D$).
  \item Columna $i$: NCube $i$ (variable futura).
  \item $T[s, i]$: costo acumulado de transitar del estado inicial al estado $s$
        evaluando el NCube $i$.
\end{itemize}
KGeoMIP accede a $T$ vía \texttt{self.\_geomip.\_flat\_T} (array C-contiguo de forma
$(2^D,\, D_{\text{nc}})$) y la reinterpreta como tensor $n$-dimensional
\texttt{\_flat\_nd} de forma $(2, 2, \ldots, 2,\, D_{\text{nc}})$ sin copiar datos
(\texttt{reshape} como vista).

\subsubsection*{Matriz de similitud causal $S$}

$S \in \mathbb{R}^{D_{\text{nc}} \times D_{\text{nc}}}$ es una matriz simétrica donde:
\begin{equation}
  S_{ij} = \frac{1}{2}\!\sum_{s=0}^{2^D - 1} T[s,\,i]
  \cdot \mathbf{1}\!\left[\text{bit}_j(s) \neq \text{bit}_j(s_0)\right]
  + \text{(término simétrico)}
  \label{eq:S}
\end{equation}
$S_{ij}$ acumula cuánto ``costo de transición'' cargó la variable futura $i$ para
llegar a estados donde la variable presente $j$ cambió respecto al estado inicial.
Cuanto mayor $S_{ij}$, más relacionadas causalmente están $i$ y $j$.
Se calcula una sola vez por subsistema (decisión D4-02) mediante multiplicación matricial
por bloques: \texttt{tabla.T @ difiere}, donde \texttt{difiere[$s$,$j$] = 1} si el bit
$j$ del estado $s$ difiere del estado inicial.

\subsubsection*{Representación de particiones}

Una k-partición se representa como \texttt{list[frozenset[int]]}: una lista de
$k$ conjuntos congelados de índices locales de NCubes. El uso de \texttt{frozenset}
es deliberado: es hashable (puede ser clave de caché y elemento de MinHeap) e inmutable
(evita mutaciones accidentales durante el refinamiento).

El MinHeap almacena tuplas donde el primer campo es el criterio de prioridad
($\Delta\Phi$ en E4, $\varphi_{\text{local}}$ en C4), garantizando extracción del
mínimo en $\mathcal{O}(\log k)$.

% ─────────────────────────────────────────────────────────────────────────────
\subsection{Pseudocódigo detallado}
% ─────────────────────────────────────────────────────────────────────────────

\subsubsection*{KGeoMIP — Algoritmo E4}

\begin{algorithm}[H]
\caption{KGeoMIP E4 — Refinamiento divisivo guiado por similitud causal}
\label{alg:kgeomip_e4}
\begin{algorithmic}[1]
\Input{Subsistema preparado (NCubes, $D$, estado inicial $s_0$), número de partes $k$}
\Output{k-partición $\Pi^* = \{P_1, \ldots, P_k\}$, pérdida $\Phi^*$}

\State \textbf{Fase 0 — Anclaje $k=2$}
\State $(\text{sol}_{k2}, T, \texttt{flat\_T}) \leftarrow \text{GeoMIP.aplicar\_estrategia}(\ldots)$
\If{$k = 2$} \Return $\text{sol}_{k2}$ \EndIf

\State \textbf{Fase 1 — Construir $S$ (una sola vez por subsistema)}
\If{$S$ no cacheada}
  \For{bloque $s$ en $\{0 \ldots 2^D-1\}$ de tamaño $\min(2^D, 2^{16})$}
    \State $\texttt{difiere}[s,j] \leftarrow \mathbf{1}[\text{bit}_j(s) \neq \text{bit}_j(s_0)]$
    \State $S \mathrel{+}= T[\text{bloque}]^\top \cdot \texttt{difiere}$
  \EndFor
  \State $S \leftarrow (S + S^\top) / 2$
\EndIf

\State \textbf{Fase 2 — Raíz: proyectar MIP de GeoMIP al espacio de NCubes}
\State $P_a \leftarrow \{\text{índice NCube } i : \text{variable}(i) \in \text{lado futuro de MIP}\}$
\State $P_b \leftarrow V \setminus P_a$
\State $(A', B', \delta') \leftarrow \text{MejorCorte}(V, S)$
\If{$\delta' < \Delta\Phi(P_a, P_b) - \varepsilon$}
  \State $P_a, P_b \leftarrow A', B'$ \Comment{raíz alternativa si es estrictamente mejor}
\EndIf
\State $\Pi \leftarrow [P_a, P_b]$

\State \textbf{Fase 3 — Refinamiento divisivo MinHeap}
\State $H \leftarrow \emptyset$ \Comment{MinHeap por $(\Delta\Phi, \texttt{id}, |P|, \min P)$}
\For{$P \in \Pi$ con $|P| \geq 2$}
  \State $(A, B, \delta) \leftarrow \text{MejorCorte}(P, S)$
  \State $\text{HeapPush}(H,\; (\delta,\, \texttt{id}^{++},\, |P|,\, \min P,\, P,\, A,\, B))$
\EndFor

\For{$t = 1$ \textbf{to} $k - |\Pi|$}
  \State $(\delta^*, \_, \_, \_, P^*, A^*, B^*) \leftarrow \text{HeapPop}(H)$
  \State $\Pi \leftarrow (\Pi \setminus \{P^*\}) \cup \{A^*\} \cup \{B^*\}$
  \For{$Q \in \{A^*, B^*\}$ con $|Q| \geq 2$}
    \State $(A_Q, B_Q, \delta_Q) \leftarrow \text{MejorCorte}(Q, S)$
    \State $\text{HeapPush}(H,\; (\delta_Q,\, \texttt{id}^{++},\, |Q|,\, \min Q,\, Q,\, A_Q,\, B_Q))$
  \EndFor
\EndFor

\State \textbf{Fase 4 — EMD-Effect real (una sola vez)}
\State $\Phi^* \leftarrow \text{CalcPhiTotal}(\Pi,\, \text{subsistema})$
\State \Return $(\Pi,\, \Phi^*)$
\end{algorithmic}
\end{algorithm}

\bigskip

\textbf{Subrutina MejorCorte($P$, $S$)} — dispatcher D4-05:

\begin{algorithm}[H]
\caption{MejorCorte — dispatcher exhaustivo / guiado por $S$}
\label{alg:mejor_corte}
\begin{algorithmic}[1]
\Input{Parte $P$ (frozenset), matriz $S$, umbral $U = 4096$}
\Output{$(A^*, B^*, \Delta\Phi^*)$}
\If{$2^{|P|-1} - 1 \leq U$}  \Comment{bloque pequeño: enumeración exacta}
  \For{cada máscara $m \in \{1, \ldots, 2^{|P|-1}-1\}$}
    \State $A \leftarrow \{P_{\text{sorted}}[i] : \text{bit}_i(m) = 1\}$,\quad $B \leftarrow P \setminus A$
    \State Actualizar $(A^*, B^*, \Delta\Phi^*)$ si $\Delta\Phi(A,B) < \Delta\Phi^*$
  \EndFor
\Else  \Comment{bloque grande: hasta 20 candidatos guiados por $S$}
  \For{$(A, B) \in$ CandidatosPorAfinidad$(P, S, 20)$}
    \State Actualizar $(A^*, B^*, \Delta\Phi^*)$ si $\Delta\Phi(A,B) < \Delta\Phi^*$
  \EndFor
\EndIf
\State \Return $(A^*, B^*, \Delta\Phi^*)$
\end{algorithmic}
\end{algorithm}

\bigskip

\subsubsection*{KQNodes — Algoritmo C4}

\begin{algorithm}[H]
\caption{KQNodes C4 — Refinamiento greedy submodular con MinHeap}
\label{alg:kqnodes_c4}
\begin{algorithmic}[1]
\Input{Subsistema preparado (NCubes, $D$, $N$, estado inicial $s_0$), número de partes $k$}
\Output{k-partición $\Pi^* = \{P_1, \ldots, P_k\}$, pérdida $\Phi^*$}

\State \textbf{Preparar tensores}
\State $\texttt{data\_nd} \leftarrow \text{np.stack}([c.\text{data}\ \forall c \in \text{NCubes}])$
  \Comment{forma $(N, 2^D)$}
\State $\texttt{pivot\_idx} \leftarrow (s_0[\text{dim}_d]\ \forall d \in \text{dims\_ncubos})$
\State $\texttt{pivot\_vals} \leftarrow \texttt{data\_nd}[\,:,\, \texttt{pivot\_idx}[::-1]]$
  \Comment{$p(\text{nodo}_i = \text{ON} \mid s_0)$}

\State $V \leftarrow \{0, 1, \ldots, \min(N,D)-1\}$
\If{$k = 2$}
  \State \Return QNodes exacto (regresión DB-03.1)
\EndIf

\State \textbf{Paso 0 — bipartir $V$ completo}
\State $(A_0, B_0, \varphi_0) \leftarrow$ QNodesBloque$(V,\, N,\, D,\, \texttt{data\_nd},\, \texttt{pivot})$
\State $H \leftarrow [(\varphi_0,\, -|V|,\, 0,\, V,\, A_0,\, B_0)]$
\State $\Pi \leftarrow [V]$,\quad $\texttt{id} \leftarrow 0$

\For{$t = 1$ \textbf{to} $k-1$}
  \State $(\varphi^*,\, \_,\, \_,\, P^*,\, A^*,\, B^*) \leftarrow \text{HeapPop}(H)$
  \State $\Pi \leftarrow (\Pi \setminus \{P^*\}) \cup \{A^*\} \cup \{B^*\}$
  \For{$Q \in \{A^*, B^*\}$ con $|Q| \geq 2$}
    \State $(A_Q, B_Q, \varphi_Q) \leftarrow$ QNodesBloque$(Q,\, N,\, D,\, \texttt{data\_nd},\, \texttt{pivot})$
    \State $\texttt{id} \mathrel{+}= 1$
    \State $\text{HeapPush}(H,\; (\varphi_Q,\, -|Q|,\, \texttt{id},\, Q,\, A_Q,\, B_Q))$
  \EndFor
\EndFor

\State \textbf{EMD-Effect real (una sola vez)}
\State $\Phi^* \leftarrow \text{CalcPhiTotal}(\Pi,\, \text{subsistema})$
\State \Return $(\Pi,\, \Phi^*)$
\end{algorithmic}
\end{algorithm}

\bigskip

\textbf{Subrutina QNodesBloque($P$)} — oracle restringido por bloque:

\begin{algorithm}[H]
\caption{QNodesBloque — bipartición de un bloque con oracle local}
\label{alg:qnodes_bloque}
\begin{algorithmic}[1]
\Input{Bloque $P \subseteq V$, tensores del subsistema, pivote}
\Output{$(A^*, B^*, \varphi_{\text{local}})$}
\State $m \leftarrow |P|$;\quad $\texttt{índices\_gl} \leftarrow \text{sorted}(P)$
\State Construir caché local $\mathcal{C} \leftarrow \{\}$ \Comment{local a esta llamada}

\Function{$f_{\text{local}}$}{máscara\_local $m_{\ell}$}
  \If{$m_{\ell} = 0$ \textbf{or} $m_{\ell} = 2^m - 1$} \Return $0$ \EndIf
  \If{$m_{\ell} \in \mathcal{C}$} \Return $\mathcal{C}[m_{\ell}].\varphi$ \EndIf
  \State Construir slices sobre \texttt{data\_nd}: libre si bit local $\in$ máscara, fijo al pivote si no
  \State $\bar{\mu}_A \leftarrow \texttt{data\_nd}[\text{slc\_A}].\text{mean}()$;\quad
         $\bar{\mu}_B \leftarrow \texttt{data\_nd}[\text{slc\_B}].\text{mean}()$
  \State $\mathcal{C}[m_{\ell}] \leftarrow (\bar{\mu}_A, \bar{\mu}_B)$;\quad
         $\mathcal{C}[\overline{m_{\ell}}] \leftarrow (\bar{\mu}_B, \bar{\mu}_A)$
         \Comment{complemento gratis}
  \State \Return $\sum_i \min(|\bar{\mu}_B[i] - v_i|,\; |\bar{\mu}_A[i] - v_i|)$
\EndFunction

\State $(\varphi^*,\, m^*) \leftarrow \text{Queyranne}(m,\, f_{\text{local}})$
  \Comment{MAO, $\mathcal{O}(m^3)$ evaluaciones}
\State $A^* \leftarrow \{\texttt{índices\_gl}[r] : \text{bit}_r(m^*) = 1\}$;\quad
       $B^* \leftarrow P \setminus A^*$
\State \Return $(A^*,\, B^*,\, \varphi^*)$
\end{algorithmic}
\end{algorithm}

% ─────────────────────────────────────────────────────────────────────────────
\subsection{Estrategia de búsqueda}
% ─────────────────────────────────────────────────────────────────────────────

\subsubsection*{KGeoMIP: candidatos guiados por afinidad causal}

Para bloques grandes ($2^{|P|-1} - 1 > 4096$), \texttt{\_mejor\_corte\_guiado\_por\_S}
genera hasta $N_{\max} = 20$ candidatos de bipartición ordenados por menor afinidad
cruzada entre A y B bajo $S$:
\begin{equation}
  \text{afinidad}(A, B) = \frac{1}{|A| \cdot |B|} \sum_{i \in A, j \in B} S_{ij}
\end{equation}
Los candidatos se seleccionan mediante \texttt{\_candidatos\_por\_afinidad} del módulo
\texttt{kgeomip\_cortes.py}, que combina tres fuentes:
\begin{enumerate}
  \item \textbf{Enumerados} (bloques $|P| \leq 5$): todas las biparticiones posibles.
  \item \textbf{Constructivos} (bloques medianos): seeds de un elemento, expandidos
        greedily por menor $S_{ij}$ hasta cubrir la mitad.
  \item \textbf{Por afinidad} (bloques grandes): cortes que separan clusters de baja
        similitud entre sí.
\end{enumerate}
De los candidatos generados, solo los primeros $N_{\max} = 20$ se evalúan con $\Delta\Phi$
real. Esto transforma una búsqueda $\mathcal{O}(2^{|P|})$ en $\mathcal{O}(|P|^2 + 20)$.

\subsubsection*{KQNodes: oracle MAO de Queyranne}

Queyranne (1998) garantiza encontrar la bipartición de mínimo de una función
submodular simétrica $f$ en exactamente $\mathcal{O}(m^3)$ evaluaciones de $f$,
donde $m = |P|$. El algoritmo construye un \textbf{Ordenamiento de Máxima Adyacencia
(MAO)}: en cada etapa agrega al ordenamiento el elemento con mayor suma de $f$ con
los ya incluidos. La bipartición de costo mínimo corresponde al último elemento del
MAO junto con el penúltimo como su «testigo».

La clave del funcionamiento correcto en KQNodes es el \textbf{caché local por bloque}:
cada llamada a \texttt{QNodesBloque} crea su propio diccionario de medias
$\mathcal{C}$. Dos bloques distintos pueden tener la misma máscara numérica local
con semántica diferente; un caché global produciría resultados incorrectos.

\subsubsection*{Evaluación de la pérdida $\Delta\Phi$ (KGeoMIP)}

La descomposición aditiva de EMD-Effect permite calcular el incremento exacto sin
recomputar $\Phi$ completo:
\begin{equation}
  \Delta\Phi(A, B) = \text{costo}(A) + \text{costo}(B) - \text{costo}(A \cup B)
  \label{eq:delta_phi}
\end{equation}
donde $\text{costo}(Q) = \sum_{d \in Q} |p_d^{\text{orig}} - p_d^{(Q)}|$
y $p_d^{(Q)}$ es la marginal del nodo $d$ condicionada solo a las dimensiones de $Q$.
Cada $\text{costo}(Q)$ se cachea en \texttt{\_costo\_cache[frozenset(Q)]}: un bloque
$P$ aparece en el heap como entero ($P$) y luego como sus hijos ($A$, $B$), de modo
que las tres evaluaciones de la Ecuación~\ref{eq:delta_phi} son hits de caché en la
segunda y tercera llamada.

% ─────────────────────────────────────────────────────────────────────────────
\subsection{Evaluación final de la k-partición}
% ─────────────────────────────────────────────────────────────────────────────

Una vez identificada la k-partición $\Pi^* = \{P_1, \ldots, P_k\}$, se calcula la
pérdida exacta mediante \textbf{CalcPhiTotal}:

\begin{enumerate}
  \item Calcular $\mathbf{p}^{\text{orig}} = \text{distribucion\_marginal}()$ del
        subsistema completo (sin partición).
  \item Para cada nodo $d$ en la parte $P_m$:
        \begin{enumerate}
          \item Obtener el NCube de $d$.
          \item Marginalizar sobre todas las dimensiones \textbf{fuera} de $P_m$
                (promedio uniforme sobre las dimensiones del mecanismo que no
                pertenecen a $P_m$).
          \item Seleccionar el valor en el estado pivote $s_0$:
                $p_d^{(\Pi)} \leftarrow \text{NCube}_d[\text{pivot}_{P_m}]$.
        \end{enumerate}
  \item $\Phi^* \leftarrow \sum_{d=1}^{N} |p_d^{\text{orig}} - p_d^{(\Pi)}|$
\end{enumerate}

Esta es la EMD-Effect bajo independencia de partes (norma $L_1$ sobre los vectores
de distribución marginal), y se evalúa \textbf{una sola vez} para la partición
ganadora. El costo es $\mathcal{O}(N \cdot D)$, despreciable frente al refinamiento.

% ─────────────────────────────────────────────────────────────────────────────
\subsection{Optimizaciones implementadas}
% ─────────────────────────────────────────────────────────────────────────────

\begin{table}[H]
  \centering
  \caption{Optimizaciones implementadas en KGeoMIP y KQNodes con su impacto.}
  \label{tab:optimizaciones}
  \begin{tabular}{p{3.2cm} p{4.5cm} p{2.5cm} p{3cm}}
    \toprule
    \textbf{Optimización} & \textbf{Descripción} & \textbf{Estrategia} & \textbf{Impacto} \\
    \midrule
    Caché de subsistema
      & Se reutilizan $S$, \texttt{flat\_nd}, sol $k=2$ y cachés entre llamadas
        con el mismo (condición, alcance, mecanismo).
      & KGeoMIP
      & Evita recalcular $S$ al barrer $k \in \{2,3,4,5\}$ \\[4pt]

    Caché de marginales \texttt{\_marg\_cache}
      & Indexado por máscara de bits entera; guarda \texttt{ndarray(D)} por evaluación.
      & KGeoMIP
      & $\mathcal{O}(1)$ para máscaras repetidas \\[4pt]

    Caché de costos \texttt{\_costo\_cache}
      & Indexado por \texttt{frozenset}; guarda el costo real de cada parte.
      & KGeoMIP
      & $\Delta\Phi$ en $\mathcal{O}(1)$ con tres lookups \\[4pt]

    Vista $n$-dimensional \texttt{\_flat\_nd}
      & \texttt{reshape} sin copia de \texttt{\_flat\_T}; permite marginalizar
        por indexación de ejes en lugar de máscaras booleanas sobre $2^D$ estados.
      & KGeoMIP
      & Marginalización en $\mathcal{O}(2^{|P|})$ en vez de $\mathcal{O}(2^D)$ \\[4pt]

    $S$ por bloques (OPT-K3)
      & Multiplicación matricial \texttt{tabla.T @ difiere} en bloques de
        $\min(2^D, 2^{16})$ filas; usa BLAS internamente.
      & KGeoMIP
      & $\sim$10$\times$ vs. bucle Python \\[4pt]

    Caché local por bloque \texttt{\_means\_cache}
      & Diccionario nuevo por llamada a \texttt{\_oracle\_restringido}; almacena
        medias $(\bar{\mu}_A, \bar{\mu}_B)$ y el complemento gratis.
      & KQNodes
      & Corte la mitad de evaluaciones del oracle \\[4pt]

    Complemento gratis del oracle
      & Al calcular $f(m_\ell)$, se guarda también $f(\overline{m_\ell})$ sin costo
        adicional (solo se invierten los roles A/B).
      & KQNodes
      & Reduce evaluaciones reales a $\sim$mitad \\[4pt]

    EMD al final (D4-04 / D3-04)
      & La EMD-Effect real se evalúa solo sobre la k-partición ganadora, no en
        cada paso del refinamiento.
      & Ambas
      & Reduce llamadas a EMD de $\mathcal{O}(kD^2)$ a $\mathcal{O}(1)$ \\[4pt]

    MinHeap \texttt{heapq}
      & Mantiene el mínimo en $\mathcal{O}(\log k)$; el heap tiene a lo sumo
        $k-1$ elementos en estado estacionario.
      & Ambas
      & $\mathcal{O}(k \log k)$ vs. $\mathcal{O}(k^2)$ por reordenamiento \\
    \bottomrule
  \end{tabular}
\end{table}
